% GENERATED FILE. Edit canonical lessons, then run npm run generate:books.
% canonical-source-hash: fnv1a64:a6aa7429a6e4fb56
% canonical-lessons: TE-C509-pirikitanam, TE-C509-amayakudu, TE-C509-murkhudu, TE-C509-baddhakam, TE-C509-husaru, TE-R509-review-502-509

\chapter{What Kind of Person}
\label{ch:te-what-kind-of-person}
\hlchaptermodality{509}

\begin{chapteropening}
\textbf{By the end of this chapter:} \emph{I can say that someone is cowardly, naive, foolish, lazy or lively.}

Complete the last lesson of chapter 509: I can say that someone is cowardly, naive, foolish, lazy or lively.
\end{chapteropening}

\section[\texorpdfstring{\te{పిరికితనం} (pirikitanaṁ)}{పిరికితనం (pirikitanaṁ)}]{\te{పిరికితనం} (pirikitanaṁ) — cowardice}

\label{lesson:TE-C509-pirikitanam}



\begin{quote}
Before the new one: say the Telugu for daring, then the Telugu for timid.
\end{quote}

\subsection*{You'll want to know: \te{పిరికితనం}}
\textbf{\te{పిరికితనం}} — \emph{pirikitanaṁ} — \textquotedblleft{}cowardice\textquotedblright{}.

\begin{quote}
\te{నిజం} \te{చెప్పకపోవడం} \te{పిరికితనం}.
\end{quote}

\begin{quote}
\emph{nijaṁ ceppakapōvaḍaṁ pirikitanaṁ.}
\end{quote}

\begin{quote}
Not telling the truth is cowardice.
\end{quote}

\te{పిరికి}, timid, + -\te{తనం}.

\begin{cousinweb}
\textbf{\te{పిరికితనం}} (\emph{pirikitanaṁ}) is native Dravidian: \emph{piṟiki}, \textquotedblleft{}timid, cowardly\textquotedblright{} (DEDR 4200), with -\te{తనం}, the ending that makes abstract nouns.
\end{cousinweb}

\begin{grammarlens}[title={What you've built}]
One more word for describing what you lived and felt.
\end{grammarlens}

\subsection*{Your turn}
\begin{itemize}
\raggedright
  \item \emph{Say it:} \emph{pirikitanaṁ}
  \item \emph{Say it:} \emph{pirikitanaṁ}, once more
  \item \emph{Say it:} \emph{nijaṁ ceppakapōvaḍaṁ pirikitanaṁ.} — Not telling the truth is cowardice.
  \item \emph{Recall:} say the Telugu for daring, then the Telugu for timid, then say \emph{pirikitanaṁ} again
\end{itemize}

\begin{tcolorbox}[breakable,colback=teal!4,colframe=teal!35!black,title={Before you move on}]
What does \te{పిరికితనం} mean? (\textquotedblleft{}Cowardice\textquotedblright{}.) Say it once more.
\end{tcolorbox}
\section[\texorpdfstring{\te{అమాయకుడు} (amāyakuḍu)}{అమాయకుడు (amāyakuḍu)}]{\te{అమాయకుడు} (amāyakuḍu) — an innocent man; a naive man}

\label{lesson:TE-C509-amayakudu}



\begin{quote}
Before the new one: say the Telugu for timid, then the Telugu for cowardice.
\end{quote}

\subsection*{You'll want to know: \te{అమాయకుడు}}
\textbf{\te{అమాయకుడు}} — \emph{amāyakuḍu} — \textquotedblleft{}an innocent man; a naive man\textquotedblright{}.

\begin{quote}
\te{ఆ} \te{అమాయకుడిని} \te{దొంగలు} \te{మోసం} \te{చేశారు}.
\end{quote}

\begin{quote}
\emph{ā amāyakuḍini doṅgalu mōsaṁ cēśāru.}
\end{quote}

\begin{quote}
Thieves cheated that naive man.
\end{quote}

Feminine: \te{అమాయకురాలు}. \te{అమాయకంగా}, innocently.

\begin{cousinweb}
\textbf{\te{అమాయకుడు}} (\emph{amāyakuḍu}) is usually traced to Sanskrit \emph{a-}, \textquotedblleft{}not\textquotedblright{}, and \emph{māyā}, which CDIAL 10060 glosses as \textquotedblleft{}illusion\textquotedblright{} and \textquotedblleft{}fraud\textquotedblright{}, with the Telugu masculine ending -\te{డు}: a man without guile.
\end{cousinweb}

\begin{grammarlens}[title={What you've built}]
One more word for describing what you lived and felt.
\end{grammarlens}

\subsection*{Your turn}
\begin{itemize}
\raggedright
  \item \emph{Say it:} \emph{amāyakuḍu}
  \item \emph{Say it:} \emph{amāyakuḍu}, once more
  \item \emph{Say it:} \emph{ā amāyakuḍini doṅgalu mōsaṁ cēśāru.} — Thieves cheated that naive man.
  \item \emph{Recall:} say the Telugu for timid, then the Telugu for cowardice, then say \emph{amāyakuḍu} again
\end{itemize}

\begin{tcolorbox}[breakable,colback=teal!4,colframe=teal!35!black,title={Before you move on}]
What does \te{అమాయకుడు} mean? (\textquotedblleft{}An innocent man; a naive man\textquotedblright{}.) Say it once more.
\end{tcolorbox}
\section[\texorpdfstring{\te{మూర్ఖుడు} (mūrkhuḍu)}{మూర్ఖుడు (mūrkhuḍu)}]{\te{మూర్ఖుడు} (mūrkhuḍu) — a fool; a stubborn fool}

\label{lesson:TE-C509-murkhudu}



\begin{quote}
Before the new one: say the Telugu for cowardice, then the Telugu for an innocent man.
\end{quote}

\subsection*{You'll want to know: \te{మూర్ఖుడు}}
\textbf{\te{మూర్ఖుడు}} — \emph{mūrkhuḍu} — \textquotedblleft{}a fool; a stubborn fool\textquotedblright{}.

\begin{quote}
\te{ఆ} \te{మూర్ఖుడు} \te{ఎవరి} \te{మాటా} \te{వినడు}.
\end{quote}

\begin{quote}
\emph{ā mūrkhuḍu evari māṭā vinaḍu.}
\end{quote}

\begin{quote}
That fool listens to no one.
\end{quote}

Someone foolish and obstinate. Feminine: \te{మూర్ఖురాలు}.

\begin{cousinweb}
\textbf{\te{మూర్ఖుడు}} (\emph{mūrkhuḍu}) is borrowed from Sanskrit \emph{mūrkha}, \textquotedblleft{}stupid, foolish; a fool\textquotedblright{} (Monier-Williams), with the Telugu masculine ending -\te{డు}.
\end{cousinweb}

\begin{grammarlens}[title={What you've built}]
One more word for describing what you lived and felt.
\end{grammarlens}

\subsection*{Your turn}
\begin{itemize}
\raggedright
  \item \emph{Say it:} \emph{mūrkhuḍu}
  \item \emph{Say it:} \emph{mūrkhuḍu}, once more
  \item \emph{Say it:} \emph{ā mūrkhuḍu evari māṭā vinaḍu.} — That fool listens to no one.
  \item \emph{Recall:} say the Telugu for cowardice, then the Telugu for an innocent man, then say \emph{mūrkhuḍu} again
\end{itemize}

\begin{tcolorbox}[breakable,colback=teal!4,colframe=teal!35!black,title={Before you move on}]
What does \te{మూర్ఖుడు} mean? (\textquotedblleft{}A fool; a stubborn fool\textquotedblright{}.) Say it once more.
\end{tcolorbox}
\section[\texorpdfstring{\te{బద్ధకం} (baddhakaṁ)}{బద్ధకం (baddhakaṁ)}]{\te{బద్ధకం} (baddhakaṁ) — laziness; sluggishness}

\label{lesson:TE-C509-baddhakam}



\begin{quote}
Before the new one: say the Telugu for an innocent man, then the Telugu for a fool.
\end{quote}

\subsection*{You'll want to know: \te{బద్ధకం}}
\textbf{\te{బద్ధకం}} — \emph{baddhakaṁ} — \textquotedblleft{}laziness; sluggishness\textquotedblright{}.

\begin{quote}
\te{ఆదివారం} \te{ఉదయం} \te{లేవడానికి} \te{నాకు} \te{బద్ధకం}.
\end{quote}

\begin{quote}
\emph{ādivāraṁ udayaṁ lēvaḍāniki nāku baddhakaṁ.}
\end{quote}

\begin{quote}
On Sunday mornings I feel too lazy to get up.
\end{quote}

\te{నాకు} \te{బద్ధకంగా} \te{ఉంది}, I feel lazy. \te{బద్ధకస్తుడు}, a lazy man.

\begin{cousinweb}
\textbf{\te{బద్ధకం}} (\emph{baddhakaṁ}) is a Telugu word for sluggishness and laziness; its origin is not settled.
\end{cousinweb}

\begin{grammarlens}[title={What you've built}]
One more word for describing what you lived and felt.
\end{grammarlens}

\subsection*{Your turn}
\begin{itemize}
\raggedright
  \item \emph{Say it:} \emph{baddhakaṁ}
  \item \emph{Say it:} \emph{baddhakaṁ}, once more
  \item \emph{Say it:} \emph{ādivāraṁ udayaṁ lēvaḍāniki nāku baddhakaṁ.} — On Sunday mornings I feel too lazy to get up.
  \item \emph{Recall:} say the Telugu for an innocent man, then the Telugu for a fool, then say \emph{baddhakaṁ} again
\end{itemize}

\begin{tcolorbox}[breakable,colback=teal!4,colframe=teal!35!black,title={Before you move on}]
What does \te{బద్ధకం} mean? (\textquotedblleft{}Laziness; sluggishness\textquotedblright{}.) Say it once more.
\end{tcolorbox}
\section[\texorpdfstring{\te{హుషారు} (huṣāru)}{హుషారు (huṣāru)}]{\te{హుషారు} (huṣāru) — high spirits; liveliness}

\label{lesson:TE-C509-husaru}



\begin{quote}
Before the new one: say the Telugu for a fool, then the Telugu for laziness.
\end{quote}

\subsection*{You'll want to know: \te{హుషారు}}
\textbf{\te{హుషారు}} — \emph{huṣāru} — \textquotedblleft{}high spirits; liveliness\textquotedblright{}.

\begin{quote}
\te{పండుగ} \te{రోజు} \te{పిల్లలు} \te{హుషారుగా} \te{ఉన్నారు}.
\end{quote}

\begin{quote}
\emph{paṇḍuga rōju pillalu huṣārugā unnāru.}
\end{quote}

\begin{quote}
On the festival day the children were full of beans.
\end{quote}

\te{హుషారుగా}, cheerfully, full of energy.

\begin{cousinweb}
\textbf{\te{హుషారు}} (\emph{huṣāru}) is usually traced to Persian \emph{hushyār}, \textquotedblleft{}alert, awake\textquotedblright{}, which came into Telugu through Hindustani; in Telugu it is the bright energy of someone in good spirits.
\end{cousinweb}

\begin{grammarlens}[title={What you've built}]
One more word for describing what you lived and felt.
\end{grammarlens}

\subsection*{Your turn}
\begin{itemize}
\raggedright
  \item \emph{Say it:} \emph{huṣāru}
  \item \emph{Say it:} \emph{huṣāru}, once more
  \item \emph{Say it:} \emph{paṇḍuga rōju pillalu huṣārugā unnāru.} — On the festival day the children were full of beans.
  \item \emph{Recall:} say the Telugu for a fool, then the Telugu for laziness, then say \emph{huṣāru} again
\end{itemize}

\begin{tcolorbox}[breakable,colback=teal!4,colframe=teal!35!black,title={Before you move on}]
What does \te{హుషారు} mean? (\textquotedblleft{}High spirits; liveliness\textquotedblright{}.) Say it once more.
\end{tcolorbox}
\section[(dialogue)]{Words from chapters 502–509 — a review}

\label{lesson:TE-R509-review-502-509}



\begin{quote}
Say the Telugu for to sigh and to be disappointed.
\end{quote}

\subsection*{Your turn}
Say these aloud:

\begin{itemize}
\raggedright
  \item to sigh, to be disappointed, discouragement, to yearn, to waste away
  \item tender love, attachment, warmth, to fondle, fondness
  \item to adore, to admire, an admirer, to appreciate, acclaim
  \item pity, compassion, to welcome warmly, to stand by, to console
  \item to tolerate, to withstand, self-restraint, peace of mind, relief
  \item temperament, personality, goodness, humility, arrogance
  \item conceit, stubbornness, determination, daring, timid
  \item cowardice, an innocent man, a fool, laziness, high spirits
\end{itemize}

\begin{tcolorbox}[breakable,colback=teal!4,colframe=teal!35!black,title={Before you move on}]
To sigh? (\textbf{\te{నిట్టూర్చు}}, \emph{niṭṭūrcu}.) To be disappointed? (\textbf{\te{నిరాశపడు}}, \emph{nirāśapaḍu}.) Discouragement? (\textbf{\te{నిరుత్సాహం}}, \emph{nirutsāhaṁ}.) To yearn? (\textbf{\te{తహతహలాడు}}, \emph{tahatahalāḍu}.) To waste away? (\textbf{\te{కృశించు}}, \emph{kṛśiñcu}.) Tender love? (\textbf{\te{అనురాగం}}, \emph{anurāgaṁ}.) Attachment? (\textbf{\te{మమకారం}}, \emph{mamakāraṁ}.) Warmth? (\textbf{\te{ఆప్యాయత}}, \emph{āpyāyata}.) To fondle? (\textbf{\te{లాలించు}}, \emph{lāliñcu}.) Fondness? (\textbf{\te{అభిమానం}}, \emph{abhimānaṁ}.) To adore? (\textbf{\te{ఆరాధించు}}, \emph{ārādhiñcu}.) To admire? (\textbf{\te{అభిమానించు}}, \emph{abhimāniñcu}.) An admirer? (\textbf{\te{అభిమాని}}, \emph{abhimāni}.) To appreciate? (\textbf{\te{మెచ్చుకో}}, \emph{meccukō}.) Acclaim? (\textbf{\te{ప్రశంస}}, \emph{praśaṁsa}.) Pity? (\textbf{\te{జాలి}}, \emph{jāli}.) Compassion? (\textbf{\te{కరుణ}}, \emph{karuṇa}.) To welcome warmly? (\textbf{\te{ఆదరించు}}, \emph{ādariñcu}.) To stand by? (\textbf{\te{ఆదుకో}}, \emph{ādukō}.) To console? (\textbf{\te{సముదాయించు}}, \emph{samudāyiñcu}.) To tolerate? (\textbf{\te{సహించు}}, \emph{sahiñcu}.) To withstand? (\textbf{\te{తట్టుకో}}, \emph{taṭṭukō}.) Self-restraint? (\textbf{\te{నిగ్రహం}}, \emph{nigrahaṁ}.) Peace of mind? (\textbf{\te{మనశ్శాంతి}}, \emph{manaśśānti}.) Relief? (\textbf{\te{ఉపశమనం}}, \emph{upaśamanaṁ}.) Temperament? (\textbf{\te{స్వభావం}}, \emph{svabhāvaṁ}.) Personality? (\textbf{\te{వ్యక్తిత్వం}}, \emph{vyaktitvaṁ}.) Goodness? (\textbf{\te{మంచితనం}}, \emph{mañcitanaṁ}.) Humility? (\textbf{\te{వినయం}}, \emph{vinayaṁ}.) Arrogance? (\textbf{\te{అహంకారం}}, \emph{ahaṅkāraṁ}.) Conceit? (\textbf{\te{పొగరు}}, \emph{pogaru}.) Stubbornness? (\textbf{\te{మొండితనం}}, \emph{moṇḍitanaṁ}.) Determination? (\textbf{\te{పట్టుదల}}, \emph{paṭṭudala}.) Daring? (\textbf{\te{సాహసం}}, \emph{sāhasaṁ}.) Timid? (\textbf{\te{పిరికి}}, \emph{piriki}.) Cowardice? (\textbf{\te{పిరికితనం}}, \emph{pirikitanaṁ}.) An innocent man? (\textbf{\te{అమాయకుడు}}, \emph{amāyakuḍu}.) A fool? (\textbf{\te{మూర్ఖుడు}}, \emph{mūrkhuḍu}.) Laziness? (\textbf{\te{బద్ధకం}}, \emph{baddhakaṁ}.) High spirits? (\textbf{\te{హుషారు}}, \emph{huṣāru}.)
\end{tcolorbox}
